%% kapa: exam=1 round=26 year=2015 period=1 form=A subject=civil_law subject_ko=민법 seq=29 number=29
%% source: Q-Net 감정평가사 기출문제 (cst00309) artlSeq=5138001 file=extracted/1008/5138001/제26회_감정평가사 1차_기출/제26회_감정평가사 1차_기출/1교시_감평_A형.pdf
%% conversion: kordoc 4.24.51 -> tools/gen1.py (auto); review=proofread-sonnet; flags=none
\begin{question}{29}
점유자와 회복자의 관계에 관한 설명으로 옳지 않은 것은? (다툼이 있으면 판례에 따름)
\begin{choices}
  \item 타주점유자가 그의 책임있는 사유로 점유물을 멸실 또는 훼손한 때에는 그가 선의로 점유했더라도 손해의 전부를 배상하여야 한다.
  \item 자신에게 과실수취권을 포함하는 권원이 있다고 오신한 데 정당한 근거가 있는 점유자는 과실수취권이 있다.
  \item 선의의 점유자가 법률상 원인 없이 회복자의 건물을 점유ㆍ사용하고 이로 말미암아 회복자에게 손해를 입혔다면 그 점유ㆍ사용으로 인한 이득을 반환할 의무가 있다.
  \item 타인 소유물을 권원 없이 점유함으로써 얻은 사용이익을 반환하는 경우, 악의의 점유자는 사용이익 뿐만 아니라 그 이자도 반환해야 한다.
  \item 점유자가 과실을 취득하였다면 통상의 필요비는 청구할 수 없다.
\end{choices}
\end{question}
\answer{3}
